% =====================================================================
%  Preámbulo compartido — Beamer
%  Curso: Análisis y Pronóstico de Datos Mineros con Python
%
%  Tema construido a mano sobre la base 'default' de Beamer:
%  compila con cualquier distribución TeX moderna sin paquetes de
%  tema externos (no requiere metropolis).
%
%  Cada deck de módulo hace:  % =====================================================================
%  Preámbulo compartido — Beamer
%  Curso: Análisis y Pronóstico de Datos Mineros con Python
%
%  Tema construido a mano sobre la base 'default' de Beamer:
%  compila con cualquier distribución TeX moderna sin paquetes de
%  tema externos (no requiere metropolis).
%
%  Cada deck de módulo hace:  % =====================================================================
%  Preámbulo compartido — Beamer
%  Curso: Análisis y Pronóstico de Datos Mineros con Python
%
%  Tema construido a mano sobre la base 'default' de Beamer:
%  compila con cualquier distribución TeX moderna sin paquetes de
%  tema externos (no requiere metropolis).
%
%  Cada deck de módulo hace:  % =====================================================================
%  Preámbulo compartido — Beamer
%  Curso: Análisis y Pronóstico de Datos Mineros con Python
%
%  Tema construido a mano sobre la base 'default' de Beamer:
%  compila con cualquier distribución TeX moderna sin paquetes de
%  tema externos (no requiere metropolis).
%
%  Cada deck de módulo hace:  \input{preambulo.tex}
% =====================================================================
\documentclass[11pt,aspectratio=169]{beamer}

\usepackage[T1]{fontenc}
\usepackage[utf8]{inputenc}
\usepackage[spanish,es-noshorthands]{babel}
\usepackage{lmodern}
\usepackage{amsmath,amssymb}
\usepackage{booktabs}
\usepackage{array}
\usepackage{xcolor}
\usepackage{graphicx}
\usepackage{listings}
\usepackage{tikz}
\usetikzlibrary{arrows.meta,positioning}

% ---------------------------------------------------------------------
%  Paleta
% ---------------------------------------------------------------------
\definecolor{azulmina}{HTML}{1F4E79}   % azul profundo
\definecolor{cobre}{HTML}{B5651D}      % acento cobre
\definecolor{griscarbon}{HTML}{2B2B2B}
\definecolor{grissuave}{HTML}{EDEEF0}
\definecolor{verdeok}{HTML}{2E7D32}
\definecolor{rojoalerta}{HTML}{C62828}

% ---------------------------------------------------------------------
%  Tema visual (sobre 'default')
% ---------------------------------------------------------------------
\usetheme{default}
\usefonttheme{professionalfonts}
\setbeamertemplate{navigation symbols}{}
\setbeamercolor{normal text}{fg=griscarbon}
\setbeamercolor{structure}{fg=azulmina}
\setbeamercolor{alerted text}{fg=rojoalerta}
\setbeamercolor{frametitle}{fg=azulmina}
\setbeamerfont{frametitle}{series=\bfseries,size=\large}
\setbeamertemplate{frametitle}{%
  \vskip6pt%
  \usebeamerfont{frametitle}\usebeamercolor[fg]{frametitle}\insertframetitle%
  \vskip2pt{\color{cobre}\hrule height 1.1pt}%
}
\setbeamertemplate{itemize items}[circle]
\setbeamertemplate{itemize subitem}[triangle]
\setbeamercolor{itemize item}{fg=cobre}
\setbeamercolor{itemize subitem}{fg=cobre}
\setbeamercolor{enumerate item}{fg=cobre}

\setbeamercolor{block title}{bg=azulmina,fg=white}
\setbeamercolor{block body}{bg=grissuave,fg=griscarbon}
\setbeamercolor{block title alerted}{bg=rojoalerta,fg=white}
\setbeamercolor{block title example}{bg=verdeok,fg=white}

% Barra inferior compacta
\setbeamertemplate{footline}{%
  \hbox{%
    \begin{beamercolorbox}[wd=.62\paperwidth,ht=2.4ex,dp=1ex,leftskip=1em]{author in head/foot}%
      \color{white}\scriptsize Análisis y Pronóstico de Datos Mineros con Python%
    \end{beamercolorbox}%
    \begin{beamercolorbox}[wd=.38\paperwidth,ht=2.4ex,dp=1ex,rightskip=1em]{title in head/foot}%
      \color{white}\scriptsize\hfill\insertshortsubtitle\ \ \textbullet\ \ \insertframenumber/\inserttotalframenumber%
    \end{beamercolorbox}}%
  \vskip0pt%
}
\setbeamercolor{author in head/foot}{bg=azulmina}
\setbeamercolor{title in head/foot}{bg=cobre}

% ---------------------------------------------------------------------
%  Código Python
% ---------------------------------------------------------------------
\lstdefinestyle{py}{
  language=Python,
  basicstyle=\ttfamily\scriptsize,
  keywordstyle=\color{azulmina}\bfseries,
  stringstyle=\color{cobre},
  commentstyle=\color{verdeok}\itshape,
  showstringspaces=false,
  breaklines=true,
  breakatwhitespace=true,
  columns=flexible,
  keepspaces=true,
  frame=leftline,
  framerule=1.4pt,
  rulecolor=\color{cobre},
  xleftmargin=1.2em,
  literate={á}{{\'a}}1 {é}{{\'e}}1 {í}{{\'i}}1 {ó}{{\'o}}1 {ú}{{\'u}}1
           {ñ}{{\~n}}1 {Á}{{\'A}}1 {É}{{\'E}}1 {Í}{{\'I}}1 {Ó}{{\'O}}1
           {Ú}{{\'U}}1 {Ñ}{{\~N}}1
}
\lstset{style=py}
\newcommand{\code}[1]{\texttt{\small #1}}

% ---------------------------------------------------------------------
%  Ayudas
% ---------------------------------------------------------------------
\newcommand{\idea}[1]{%
  \begin{center}\setlength{\fboxsep}{8pt}%
  \colorbox{grissuave}{\parbox{0.9\linewidth}{\centering\color{azulmina}\bfseries #1}}%
  \end{center}}
\newcommand{\ojo}[1]{{\color{rojoalerta}\textbf{Ojo:} #1}}
\newenvironment{recordar}{\begin{block}{Qué deberías recordar}}{\end{block}}

% flecha de flujo (segura en modo texto y modo matemático)
\newcommand{\fl}{\ensuremath{\;{\color{cobre}\rightarrow}\;}}

% cadena de flujo (roadmap): se escala para caber en el ancho de la diapositiva
\newcommand{\ruta}[1]{\par\vskip2pt\begin{center}%
  \resizebox{0.97\linewidth}{!}{\ensuremath{\displaystyle #1}}%
  \end{center}\vskip2pt}

% ---------------------------------------------------------------------
%  Metadatos
% ---------------------------------------------------------------------
\title[Datos Mineros con Python]{Análisis y Pronóstico de Datos Mineros con Python}
\author{}
\date{16 h \textbullet\ 8 clases \textbullet\ Google Colab}

% Portada estándar del módulo:  \portada{N}{Título del módulo}
\newcommand{\portada}[2]{%
  \subtitle[Módulo #1]{Módulo #1 \textemdash\ #2}%
  \begin{frame}[plain]
    \vfill
    {\color{azulmina}\rule{\linewidth}{2pt}}\\[0.4em]
    {\large\bfseries\color{cobre}Módulo #1}\\[0.3em]
    {\LARGE\bfseries\color{azulmina} #2}\\[0.6em]
    {\color{azulmina}\rule{\linewidth}{2pt}}\\[1.2em]
    \textit{Análisis y Pronóstico de Datos Mineros con Python}\\[0.2em]
    \footnotesize Apuntes + notebook de Google Colab
    \vfill
  \end{frame}}

% =====================================================================
\documentclass[11pt,aspectratio=169]{beamer}

\usepackage[T1]{fontenc}
\usepackage[utf8]{inputenc}
\usepackage[spanish,es-noshorthands]{babel}
\usepackage{lmodern}
\usepackage{amsmath,amssymb}
\usepackage{booktabs}
\usepackage{array}
\usepackage{xcolor}
\usepackage{graphicx}
\usepackage{listings}
\usepackage{tikz}
\usetikzlibrary{arrows.meta,positioning}

% ---------------------------------------------------------------------
%  Paleta
% ---------------------------------------------------------------------
\definecolor{azulmina}{HTML}{1F4E79}   % azul profundo
\definecolor{cobre}{HTML}{B5651D}      % acento cobre
\definecolor{griscarbon}{HTML}{2B2B2B}
\definecolor{grissuave}{HTML}{EDEEF0}
\definecolor{verdeok}{HTML}{2E7D32}
\definecolor{rojoalerta}{HTML}{C62828}

% ---------------------------------------------------------------------
%  Tema visual (sobre 'default')
% ---------------------------------------------------------------------
\usetheme{default}
\usefonttheme{professionalfonts}
\setbeamertemplate{navigation symbols}{}
\setbeamercolor{normal text}{fg=griscarbon}
\setbeamercolor{structure}{fg=azulmina}
\setbeamercolor{alerted text}{fg=rojoalerta}
\setbeamercolor{frametitle}{fg=azulmina}
\setbeamerfont{frametitle}{series=\bfseries,size=\large}
\setbeamertemplate{frametitle}{%
  \vskip6pt%
  \usebeamerfont{frametitle}\usebeamercolor[fg]{frametitle}\insertframetitle%
  \vskip2pt{\color{cobre}\hrule height 1.1pt}%
}
\setbeamertemplate{itemize items}[circle]
\setbeamertemplate{itemize subitem}[triangle]
\setbeamercolor{itemize item}{fg=cobre}
\setbeamercolor{itemize subitem}{fg=cobre}
\setbeamercolor{enumerate item}{fg=cobre}

\setbeamercolor{block title}{bg=azulmina,fg=white}
\setbeamercolor{block body}{bg=grissuave,fg=griscarbon}
\setbeamercolor{block title alerted}{bg=rojoalerta,fg=white}
\setbeamercolor{block title example}{bg=verdeok,fg=white}

% Barra inferior compacta
\setbeamertemplate{footline}{%
  \hbox{%
    \begin{beamercolorbox}[wd=.62\paperwidth,ht=2.4ex,dp=1ex,leftskip=1em]{author in head/foot}%
      \color{white}\scriptsize Análisis y Pronóstico de Datos Mineros con Python%
    \end{beamercolorbox}%
    \begin{beamercolorbox}[wd=.38\paperwidth,ht=2.4ex,dp=1ex,rightskip=1em]{title in head/foot}%
      \color{white}\scriptsize\hfill\insertshortsubtitle\ \ \textbullet\ \ \insertframenumber/\inserttotalframenumber%
    \end{beamercolorbox}}%
  \vskip0pt%
}
\setbeamercolor{author in head/foot}{bg=azulmina}
\setbeamercolor{title in head/foot}{bg=cobre}

% ---------------------------------------------------------------------
%  Código Python
% ---------------------------------------------------------------------
\lstdefinestyle{py}{
  language=Python,
  basicstyle=\ttfamily\scriptsize,
  keywordstyle=\color{azulmina}\bfseries,
  stringstyle=\color{cobre},
  commentstyle=\color{verdeok}\itshape,
  showstringspaces=false,
  breaklines=true,
  breakatwhitespace=true,
  columns=flexible,
  keepspaces=true,
  frame=leftline,
  framerule=1.4pt,
  rulecolor=\color{cobre},
  xleftmargin=1.2em,
  literate={á}{{\'a}}1 {é}{{\'e}}1 {í}{{\'i}}1 {ó}{{\'o}}1 {ú}{{\'u}}1
           {ñ}{{\~n}}1 {Á}{{\'A}}1 {É}{{\'E}}1 {Í}{{\'I}}1 {Ó}{{\'O}}1
           {Ú}{{\'U}}1 {Ñ}{{\~N}}1
}
\lstset{style=py}
\newcommand{\code}[1]{\texttt{\small #1}}

% ---------------------------------------------------------------------
%  Ayudas
% ---------------------------------------------------------------------
\newcommand{\idea}[1]{%
  \begin{center}\setlength{\fboxsep}{8pt}%
  \colorbox{grissuave}{\parbox{0.9\linewidth}{\centering\color{azulmina}\bfseries #1}}%
  \end{center}}
\newcommand{\ojo}[1]{{\color{rojoalerta}\textbf{Ojo:} #1}}
\newenvironment{recordar}{\begin{block}{Qué deberías recordar}}{\end{block}}

% flecha de flujo (segura en modo texto y modo matemático)
\newcommand{\fl}{\ensuremath{\;{\color{cobre}\rightarrow}\;}}

% cadena de flujo (roadmap): se escala para caber en el ancho de la diapositiva
\newcommand{\ruta}[1]{\par\vskip2pt\begin{center}%
  \resizebox{0.97\linewidth}{!}{\ensuremath{\displaystyle #1}}%
  \end{center}\vskip2pt}

% ---------------------------------------------------------------------
%  Metadatos
% ---------------------------------------------------------------------
\title[Datos Mineros con Python]{Análisis y Pronóstico de Datos Mineros con Python}
\author{}
\date{16 h \textbullet\ 8 clases \textbullet\ Google Colab}

% Portada estándar del módulo:  \portada{N}{Título del módulo}
\newcommand{\portada}[2]{%
  \subtitle[Módulo #1]{Módulo #1 \textemdash\ #2}%
  \begin{frame}[plain]
    \vfill
    {\color{azulmina}\rule{\linewidth}{2pt}}\\[0.4em]
    {\large\bfseries\color{cobre}Módulo #1}\\[0.3em]
    {\LARGE\bfseries\color{azulmina} #2}\\[0.6em]
    {\color{azulmina}\rule{\linewidth}{2pt}}\\[1.2em]
    \textit{Análisis y Pronóstico de Datos Mineros con Python}\\[0.2em]
    \footnotesize Apuntes + notebook de Google Colab
    \vfill
  \end{frame}}

% =====================================================================
\documentclass[11pt,aspectratio=169]{beamer}

\usepackage[T1]{fontenc}
\usepackage[utf8]{inputenc}
\usepackage[spanish,es-noshorthands]{babel}
\usepackage{lmodern}
\usepackage{amsmath,amssymb}
\usepackage{booktabs}
\usepackage{array}
\usepackage{xcolor}
\usepackage{graphicx}
\usepackage{listings}
\usepackage{tikz}
\usetikzlibrary{arrows.meta,positioning}

% ---------------------------------------------------------------------
%  Paleta
% ---------------------------------------------------------------------
\definecolor{azulmina}{HTML}{1F4E79}   % azul profundo
\definecolor{cobre}{HTML}{B5651D}      % acento cobre
\definecolor{griscarbon}{HTML}{2B2B2B}
\definecolor{grissuave}{HTML}{EDEEF0}
\definecolor{verdeok}{HTML}{2E7D32}
\definecolor{rojoalerta}{HTML}{C62828}

% ---------------------------------------------------------------------
%  Tema visual (sobre 'default')
% ---------------------------------------------------------------------
\usetheme{default}
\usefonttheme{professionalfonts}
\setbeamertemplate{navigation symbols}{}
\setbeamercolor{normal text}{fg=griscarbon}
\setbeamercolor{structure}{fg=azulmina}
\setbeamercolor{alerted text}{fg=rojoalerta}
\setbeamercolor{frametitle}{fg=azulmina}
\setbeamerfont{frametitle}{series=\bfseries,size=\large}
\setbeamertemplate{frametitle}{%
  \vskip6pt%
  \usebeamerfont{frametitle}\usebeamercolor[fg]{frametitle}\insertframetitle%
  \vskip2pt{\color{cobre}\hrule height 1.1pt}%
}
\setbeamertemplate{itemize items}[circle]
\setbeamertemplate{itemize subitem}[triangle]
\setbeamercolor{itemize item}{fg=cobre}
\setbeamercolor{itemize subitem}{fg=cobre}
\setbeamercolor{enumerate item}{fg=cobre}

\setbeamercolor{block title}{bg=azulmina,fg=white}
\setbeamercolor{block body}{bg=grissuave,fg=griscarbon}
\setbeamercolor{block title alerted}{bg=rojoalerta,fg=white}
\setbeamercolor{block title example}{bg=verdeok,fg=white}

% Barra inferior compacta
\setbeamertemplate{footline}{%
  \hbox{%
    \begin{beamercolorbox}[wd=.62\paperwidth,ht=2.4ex,dp=1ex,leftskip=1em]{author in head/foot}%
      \color{white}\scriptsize Análisis y Pronóstico de Datos Mineros con Python%
    \end{beamercolorbox}%
    \begin{beamercolorbox}[wd=.38\paperwidth,ht=2.4ex,dp=1ex,rightskip=1em]{title in head/foot}%
      \color{white}\scriptsize\hfill\insertshortsubtitle\ \ \textbullet\ \ \insertframenumber/\inserttotalframenumber%
    \end{beamercolorbox}}%
  \vskip0pt%
}
\setbeamercolor{author in head/foot}{bg=azulmina}
\setbeamercolor{title in head/foot}{bg=cobre}

% ---------------------------------------------------------------------
%  Código Python
% ---------------------------------------------------------------------
\lstdefinestyle{py}{
  language=Python,
  basicstyle=\ttfamily\scriptsize,
  keywordstyle=\color{azulmina}\bfseries,
  stringstyle=\color{cobre},
  commentstyle=\color{verdeok}\itshape,
  showstringspaces=false,
  breaklines=true,
  breakatwhitespace=true,
  columns=flexible,
  keepspaces=true,
  frame=leftline,
  framerule=1.4pt,
  rulecolor=\color{cobre},
  xleftmargin=1.2em,
  literate={á}{{\'a}}1 {é}{{\'e}}1 {í}{{\'i}}1 {ó}{{\'o}}1 {ú}{{\'u}}1
           {ñ}{{\~n}}1 {Á}{{\'A}}1 {É}{{\'E}}1 {Í}{{\'I}}1 {Ó}{{\'O}}1
           {Ú}{{\'U}}1 {Ñ}{{\~N}}1
}
\lstset{style=py}
\newcommand{\code}[1]{\texttt{\small #1}}

% ---------------------------------------------------------------------
%  Ayudas
% ---------------------------------------------------------------------
\newcommand{\idea}[1]{%
  \begin{center}\setlength{\fboxsep}{8pt}%
  \colorbox{grissuave}{\parbox{0.9\linewidth}{\centering\color{azulmina}\bfseries #1}}%
  \end{center}}
\newcommand{\ojo}[1]{{\color{rojoalerta}\textbf{Ojo:} #1}}
\newenvironment{recordar}{\begin{block}{Qué deberías recordar}}{\end{block}}

% flecha de flujo (segura en modo texto y modo matemático)
\newcommand{\fl}{\ensuremath{\;{\color{cobre}\rightarrow}\;}}

% cadena de flujo (roadmap): se escala para caber en el ancho de la diapositiva
\newcommand{\ruta}[1]{\par\vskip2pt\begin{center}%
  \resizebox{0.97\linewidth}{!}{\ensuremath{\displaystyle #1}}%
  \end{center}\vskip2pt}

% ---------------------------------------------------------------------
%  Metadatos
% ---------------------------------------------------------------------
\title[Datos Mineros con Python]{Análisis y Pronóstico de Datos Mineros con Python}
\author{}
\date{16 h \textbullet\ 8 clases \textbullet\ Google Colab}

% Portada estándar del módulo:  \portada{N}{Título del módulo}
\newcommand{\portada}[2]{%
  \subtitle[Módulo #1]{Módulo #1 \textemdash\ #2}%
  \begin{frame}[plain]
    \vfill
    {\color{azulmina}\rule{\linewidth}{2pt}}\\[0.4em]
    {\large\bfseries\color{cobre}Módulo #1}\\[0.3em]
    {\LARGE\bfseries\color{azulmina} #2}\\[0.6em]
    {\color{azulmina}\rule{\linewidth}{2pt}}\\[1.2em]
    \textit{Análisis y Pronóstico de Datos Mineros con Python}\\[0.2em]
    \footnotesize Apuntes + notebook de Google Colab
    \vfill
  \end{frame}}

% =====================================================================
\documentclass[11pt,aspectratio=169]{beamer}

\usepackage[T1]{fontenc}
\usepackage[utf8]{inputenc}
\usepackage[spanish,es-noshorthands]{babel}
\usepackage{lmodern}
\usepackage{amsmath,amssymb}
\usepackage{booktabs}
\usepackage{array}
\usepackage{xcolor}
\usepackage{graphicx}
\usepackage{listings}
\usepackage{tikz}
\usetikzlibrary{arrows.meta,positioning}

% ---------------------------------------------------------------------
%  Paleta
% ---------------------------------------------------------------------
\definecolor{azulmina}{HTML}{1F4E79}   % azul profundo
\definecolor{cobre}{HTML}{B5651D}      % acento cobre
\definecolor{griscarbon}{HTML}{2B2B2B}
\definecolor{grissuave}{HTML}{EDEEF0}
\definecolor{verdeok}{HTML}{2E7D32}
\definecolor{rojoalerta}{HTML}{C62828}

% ---------------------------------------------------------------------
%  Tema visual (sobre 'default')
% ---------------------------------------------------------------------
\usetheme{default}
\usefonttheme{professionalfonts}
\setbeamertemplate{navigation symbols}{}
\setbeamercolor{normal text}{fg=griscarbon}
\setbeamercolor{structure}{fg=azulmina}
\setbeamercolor{alerted text}{fg=rojoalerta}
\setbeamercolor{frametitle}{fg=azulmina}
\setbeamerfont{frametitle}{series=\bfseries,size=\large}
\setbeamertemplate{frametitle}{%
  \vskip6pt%
  \usebeamerfont{frametitle}\usebeamercolor[fg]{frametitle}\insertframetitle%
  \vskip2pt{\color{cobre}\hrule height 1.1pt}%
}
\setbeamertemplate{itemize items}[circle]
\setbeamertemplate{itemize subitem}[triangle]
\setbeamercolor{itemize item}{fg=cobre}
\setbeamercolor{itemize subitem}{fg=cobre}
\setbeamercolor{enumerate item}{fg=cobre}

\setbeamercolor{block title}{bg=azulmina,fg=white}
\setbeamercolor{block body}{bg=grissuave,fg=griscarbon}
\setbeamercolor{block title alerted}{bg=rojoalerta,fg=white}
\setbeamercolor{block title example}{bg=verdeok,fg=white}

% Barra inferior compacta
\setbeamertemplate{footline}{%
  \hbox{%
    \begin{beamercolorbox}[wd=.62\paperwidth,ht=2.4ex,dp=1ex,leftskip=1em]{author in head/foot}%
      \color{white}\scriptsize Análisis y Pronóstico de Datos Mineros con Python%
    \end{beamercolorbox}%
    \begin{beamercolorbox}[wd=.38\paperwidth,ht=2.4ex,dp=1ex,rightskip=1em]{title in head/foot}%
      \color{white}\scriptsize\hfill\insertshortsubtitle\ \ \textbullet\ \ \insertframenumber/\inserttotalframenumber%
    \end{beamercolorbox}}%
  \vskip0pt%
}
\setbeamercolor{author in head/foot}{bg=azulmina}
\setbeamercolor{title in head/foot}{bg=cobre}

% ---------------------------------------------------------------------
%  Código Python
% ---------------------------------------------------------------------
\lstdefinestyle{py}{
  language=Python,
  basicstyle=\ttfamily\scriptsize,
  keywordstyle=\color{azulmina}\bfseries,
  stringstyle=\color{cobre},
  commentstyle=\color{verdeok}\itshape,
  showstringspaces=false,
  breaklines=true,
  breakatwhitespace=true,
  columns=flexible,
  keepspaces=true,
  frame=leftline,
  framerule=1.4pt,
  rulecolor=\color{cobre},
  xleftmargin=1.2em,
  literate={á}{{\'a}}1 {é}{{\'e}}1 {í}{{\'i}}1 {ó}{{\'o}}1 {ú}{{\'u}}1
           {ñ}{{\~n}}1 {Á}{{\'A}}1 {É}{{\'E}}1 {Í}{{\'I}}1 {Ó}{{\'O}}1
           {Ú}{{\'U}}1 {Ñ}{{\~N}}1
}
\lstset{style=py}
\newcommand{\code}[1]{\texttt{\small #1}}

% ---------------------------------------------------------------------
%  Ayudas
% ---------------------------------------------------------------------
\newcommand{\idea}[1]{%
  \begin{center}\setlength{\fboxsep}{8pt}%
  \colorbox{grissuave}{\parbox{0.9\linewidth}{\centering\color{azulmina}\bfseries #1}}%
  \end{center}}
\newcommand{\ojo}[1]{{\color{rojoalerta}\textbf{Ojo:} #1}}
\newenvironment{recordar}{\begin{block}{Qué deberías recordar}}{\end{block}}

% flecha de flujo (segura en modo texto y modo matemático)
\newcommand{\fl}{\ensuremath{\;{\color{cobre}\rightarrow}\;}}

% cadena de flujo (roadmap): se escala para caber en el ancho de la diapositiva
\newcommand{\ruta}[1]{\par\vskip2pt\begin{center}%
  \resizebox{0.97\linewidth}{!}{\ensuremath{\displaystyle #1}}%
  \end{center}\vskip2pt}

% ---------------------------------------------------------------------
%  Metadatos
% ---------------------------------------------------------------------
\title[Datos Mineros con Python]{Análisis y Pronóstico de Datos Mineros con Python}
\author{}
\date{16 h \textbullet\ 8 clases \textbullet\ Google Colab}

% Portada estándar del módulo:  \portada{N}{Título del módulo}
\newcommand{\portada}[2]{%
  \subtitle[Módulo #1]{Módulo #1 \textemdash\ #2}%
  \begin{frame}[plain]
    \vfill
    {\color{azulmina}\rule{\linewidth}{2pt}}\\[0.4em]
    {\large\bfseries\color{cobre}Módulo #1}\\[0.3em]
    {\LARGE\bfseries\color{azulmina} #2}\\[0.6em]
    {\color{azulmina}\rule{\linewidth}{2pt}}\\[1.2em]
    \textit{Análisis y Pronóstico de Datos Mineros con Python}\\[0.2em]
    \footnotesize Apuntes + notebook de Google Colab
    \vfill
  \end{frame}}
